\subsection{Derinimo ir stebėjimo modelis}
\label{sec:derinimo-modelis}

Kadangi emuliatorius paleidžia pilną Linux aplinką, klaidų paieška vien pagal terminalo išvestį būtų nepakankama. Dėl to projektuotas nuotolinio derinimo mechanizmas, suderinamas su GDB Remote Serial Protocol \cite{gdb-remote-protocol}. Derinimo serveris veikia per UNIX lizdą, o procesoriaus vykdymo ciklas kiekvieno žingsnio pradžioje patikrina, ar GDB serveris leidžia tęsti vykdymą.

Derinimo sąsaja projektuota ne kaip atskiras procesoriaus kopijavimas, bet kaip siauras tarpinis sluoksnis tarp GDB serverio ir emuliuojamos sistemos. Jis leidžia GDB serveriui nuskaityti bendrosios paskirties registrus, programos skaitiklį, atminties sritis ir kelias virtualios atminties diagnostikos reikšmes. Tokiu būdu GDB serveris nepriklauso nuo vidinės procesoriaus realizacijos daugiau negu būtina.

Be standartinių GDB operacijų, numatyti keli darbo metu naudingi monitor tipo veiksmai:

\begin{itemize}
    \item \texttt{monitor mempages} -- išveda aktyvios arba nurodytos puslapių lentelės įrašus;
    \item \texttt{monitor translate <addr>} -- bando išversti virtualų adresą į fizinį;
\end{itemize}

Šios komandos nėra RISC-V architektūros dalis. Jos buvo projektuotos kaip praktinė priemonė patikrinti, ar Linux korektiškai sukuria puslapių lenteles, ar \texttt{satp} būsena atitinka laukiamą adresų erdvę ir ar išimtys kyla dėl instrukcijos dekodavimo, ar dėl adresų vertimo klaidos.
